
\textbf{Proposition 2.3.1.} \uline{Pour tout schéma propre} @@M0@@ \uline{sur} @@M1@@\uline{, et pour tout} @@M2@@\uline{, l'application}
@@M3@@
\uline{définie par}
@@M4@@
\uline{est une application polynôme de degré} @@M5@@.

Lorsque @@M6@@ sont des @@M7@@-modules inversibles et @@M8@@ un @@M9@@-module cohérent tel que @@M10@@, alors 2.3.1. implique en particulier que la caractéristique d'Euler-Poincaré
@@M11@@
est un polynôme numérique dans les variables @@M12@@, de degré @@M13@@, parfois appelé \emph{polynôme de Snapper} pour @@M14@@ et la suite @@M15@@.

\subsection*{2.4.}
Rappelons [2] que le foncteur de Picard @@M16@@ de @@M17@@ sur @@M18@@ est le faisceau en groupes dans la topologie (@@M19@@) de @@M20@@ associé au préfaisceau
@@M21@@
Lorsque @@M22@@ est algébriquement clos, on sait que l'homomorphisme canonique
@@M23@@
est un isomorphisme.



% Local continuation macros
\providecommand{\Filt}{\operatorname{Fil}}
\providecommand{\fpqc}{\mathrm{fpqc}}
\providecommand{\Sch}{\operatorname{Sch}}
\renewcommand{\Ob}{\operatorname{ob}}
\providecommand{\Gm}{\mathbb G_m}
\providecommand{\Der}{\operatorname{Der}}
\providecommand{\SL}{\operatorname{SL}}
\providecommand{\Aut}{\operatorname{Aut}}
\providecommand{\Mod}{\operatorname{Mod}}
\providecommand{\Tors}{\operatorname{Tors}}
\providecommand{\Mat}{\operatorname{Mat}}
\providecommand{\len}{\operatorname{long}}
\renewcommand{\id}{\operatorname{id}}
\providecommand{\Spec}{\operatorname{Spec}}
\providecommand{\Picf}{\operatorname{Pic}}
\providecommand{\Ens}{\operatorname{Ens}}
\providecommand{\clp}{\operatorname{cl}_{P}}
\providecommand{\clo}{\operatorname{cl}^{\circ}}
\providecommand{\clb}{\operatorname{cl}_{\bullet}}
\providecommand{\Grb}{\operatorname{Gr}_{\bullet}}
\providecommand{\Gro}{\operatorname{Gr}^{\bullet}}
\providecommand{\longO}[1]
